\subsection*{集中脉冲与影响函数}
\phantomsection\label{sec:green-response}
定理~\ref{thm:常数变易公式} 已把非齐次解写成初值演化与外部作用的积分。下面从标量情形 $y'=a(t)y+f(t)$ 出发，说明每个时刻的外部作用
如何影响以后时刻的状态。为分离这些作用，先研究集中在一个时刻的扰动。

\begin{definition}[$\delta$ 型序列]
设 $h_N\geq0$ 为光滑函数，$\displaystyle\int_\R h_N(s)\,\dd s=1$，
且存在 $\varepsilon_N\downarrow0$，使 $h_N$ 在 $[-\varepsilon_N,\varepsilon_N]$ 外为零。
称 $(h_N)$ 为一个 $\delta$ 型序列。记号 $\delta$ 表示其对测试函数积分的极限，
不是一个处处有限的通常函数。
\end{definition}
\begin{lemma}
若 $f$ 在 $\xi$ 附近连续，则
\[
 \lim_{N\to\infty}\int_\R h_N(s-\xi)f(s)\,\dd s=f(\xi).
\]
\end{lemma}
\begin{proof}
因为脉冲非负且积分为 1，先把常数 $f(\xi)$ 写成
$f(\xi)\displaystyle\int h_N(s-\xi)\,ds$。于是
\begin{align*}
 \left|\int h_N(s-\xi)f(s)\,ds-f(\xi)\right|
 &=\left|\int h_N(s-\xi)(f(s)-f(\xi))\,ds\right|\\
 &\leq\sup_{|s-\xi|\leq\varepsilon_N}|f(s)-f(\xi)|
       \int h_N(s-\xi)\,ds\\
 &=\sup_{|s-\xi|\leq\varepsilon_N}|f(s)-f(\xi)|.
\end{align*}
这里用支集条件把所有非零贡献限制在 $\xi$ 的小邻域内。
给定任意 $\eta>0$，连续性给出 $\delta>0$，使 $|s-\xi|<\delta$ 时
$|f(s)-f(\xi)|<\eta$。充分大的 $N$ 满足 $\varepsilon_N<\delta$，
所以上述差小于 $\eta$，极限成立。
\end{proof}

\begin{definition}[推迟影响函数]
对 $y'=a(t)y+f(t)$，给定零初值 $y(t_0)=0$。
在 $\xi>t_0$ 加入单位集中脉冲所产生的极限响应，称为在 $\xi$ 施加扰动的
\textbf{影响函数}，亦称推迟 Green 函数，记作 $G(t,\xi)$。
“推迟”表示扰动发生之前响应为零。
\end{definition}
\begin{theorem}
设 $a$ 连续，则在 $t\ne\xi$ 时
\[
 G(t,\xi)=
 \begin{cases}
 0,&t<\xi,\\
 \exp\displaystyle\int_\xi^t a(s)\,\dd s,&t>\xi.
 \end{cases}
\]
在 $t=\xi$ 处的右极限与左极限之差为 1。对连续 $f$ 和 $t\geq t_0$，
\[
 y(t)=\int_{t_0}^tG(t,\xi)f(\xi)\,\dd\xi
\]
是零初值问题的唯一解。
\end{theorem}
\begin{proof}
把 $f(t)$ 换成 $h_N(t-\xi)$。由本章已证的常数变易公式，响应为
\[
 y_N(t)=\int_{t_0}^t
 \exp\left(\int_s^t a(u)\,\dd u\right)h_N(s-\xi)\,\dd s.
\]
固定 $t<\xi$ 时，充分大的 $N$ 使积分区间与脉冲支集不相交，因此极限为零。
固定 $t>\xi$ 时，脉冲支集最终包含在积分区间内部，前一引理给出所述指数函数。
由 $a$ 连续，指数在 $t\downarrow\xi$ 时趋于 1，所以跳跃为 1。

最后的积分公式正是原有常数变易公式。也可对它求导：变动上限产生 $f(t)$，
核对 $t$ 的导数产生 $a(t)y(t)$，在 $t=t_0$ 积分为零。唯一性由
定理~\ref{thm:线性方程组的存在唯一性} 保证。
\end{proof}
\begin{remark}
边界点 $t=\xi$ 处的单点赋值不影响上式积分。
如果脉冲恰好位于初值时刻，应另外约定左右侧初值，不能把脉冲跨越的跳跃忽略。
向量情形把指数核换成状态转移矩阵 $X(t)X(s)^{-1}$，证明不变。
\end{remark}

图~\ref{fig:ode-tikz-pulse-green} 将光滑脉冲的集中过程与极限响应分开表示；空心点不指定 $t=\xi$ 处的单点取值。
\begin{center}
\begin{minipage}{0.96\linewidth}
\centering
% TikZ mathematical illustration; compiled with the main document.
\begin{tikzpicture}[>=Stealth,font=\small]
\begin{scope}[x=1.1cm,y=1cm]
 \draw[->,gray] (0,0)--(4,0) node[right]{$t$};
 \draw[->,gray] (.2,0)--(.2,3.2) node[above]{$h_N(t-\xi)$};
 \draw[gray,dashed] (2,0)--(2,3);
\draw[GreenMain,thick,domain=1.00100:2.99900,samples=90] plot (\x,{exp(1-1/(1-((\x-2)/1.0)^2))/(1.0*1.20690032)});
\draw[GoldMain,thick,domain=1.50050:2.49950,samples=90] plot (\x,{exp(1-1/(1-((\x-2)/0.5)^2))/(0.5*1.20690032)});
\draw[GreenDark,thick,domain=1.70030:2.29970,samples=90] plot (\x,{exp(1-1/(1-((\x-2)/0.3)^2))/(0.3*1.20690032)});

 \node[below] at (2,0){$\xi$};\node at (2,-.6){支集缩小，面积保持为 1};
\end{scope}
\begin{scope}[shift={(6.5,0)},x=1.1cm,y=1.8cm]
 \draw[->,gray] (0,0)--(4.4,0) node[right]{$t$};
 \draw[->,gray] (0,0)--(0,1.7) node[above]{$G(t,\xi)$};
 \draw[GreenMain,thick] (0,0)--(1,0);
 \draw[GreenMain,thick,domain=1:4,samples=65] plot (\x,{exp(-.7*(\x-1))});
 \draw[GoldMain,dashed] (1,0)--(1,1);
 \draw[fill=white,GreenMain] (1,0) circle(1.8pt);\draw[fill=white,GreenMain] (1,1) circle(1.8pt);
 \node[below] at (1,0){$\xi$};\node[left] at (0,1){$1$};
 \node at (2,-.34){$a(t)=-0.7$};
\end{scope}
\end{tikzpicture}

\captionof{figure}{左：单位面积光滑集中脉冲。右：$a=-0.7$ 时推迟影响函数在扰动时刻产生单位跳跃。}
\label{fig:ode-tikz-pulse-green}
\end{minipage}
\end{center}
